\documentclass[11pt,reqno]{amsart}
\usepackage[T1]{fontenc}
\usepackage[utf8]{inputenc}
\usepackage{amsmath,amssymb,amsthm,mathtools}
\usepackage{microtype}\microtypesetup{expansion=false}
\usepackage{hyperref}
\usepackage{url}
\usepackage{booktabs}
\usepackage{caption}
\usepackage{enumitem}
\usepackage{xcolor}
\usepackage{graphicx}

\hypersetup{
  colorlinks=true,
  linkcolor=blue,
  citecolor=blue,
  urlcolor=blue
}

\newtheorem{theorem}{Theorem}[section]
\newtheorem{proposition}[theorem]{Proposition}
\newtheorem{lemma}[theorem]{Lemma}
\newtheorem{corollary}[theorem]{Corollary}
\newtheorem{conjecture}[theorem]{Conjecture}
\theoremstyle{definition}
\newtheorem{definition}[theorem]{Definition}
\newtheorem{remark}[theorem]{Remark}

\newcommand{\midset}{\mathrm{mid}}
\newcommand{\BOTTOM}{\mathrm{BOTTOM}}
\newcommand{\RIGHT}{\mathrm{RIGHT}}
\newcommand{\TOP}{\mathrm{TOP}}
\newcommand{\LEFT}{\mathrm{LEFT}}

\title[Six unit squares cannot cover a larger square]
{Six unit squares cannot cover a square of side greater than~$2$}

\author{AI assistant (released by Ningyu Zhou)}
\address{Public GitHub release maintained by Ningyu Zhou, School of Physics, Peking University}
\email{zhouningyu27@gmail.com}
\urladdr{https://github.com/ZhouNingyu327/S6-eq-2}
\thanks{This manuscript was produced by a generative AI system as a \emph{candidate}
computer-assisted argument. The human releaser hosts the package and does not claim
independent mastery of every step. Independent peer review has not occurred.
See the repository \texttt{DISCLOSURE.md}.}

\date{October 21, 2026}

\begin{document}
\begin{abstract}
Let $S(n)$ be the largest side length of a Euclidean square that can be covered
by $n$ closed unit squares, allowing arbitrary rigid motions and overlaps.
Equivalently, Soifer's function $\Pi(2)$ is the least number of unit squares
needed to cover \emph{some} square of side $2+\varepsilon$ for \emph{some}
$\varepsilon>0$.
We present a computer-assisted \emph{candidate} argument that $S(6)=2$, or
equivalently $\Pi(2)=7$, addressing the conjecture stated by
D\'osa--L\'angi--Tuza~\cite{DLT2026}.
In particular, the candidate claims that six unit squares cannot cover any
square of side strictly greater than~$2$, while side~$2$ is attained by the
obvious $2\times 2$ tiling.  The matching upper bound $\Pi(2)\le 7$ is the
classical fact that seven unit squares cover some square of side $>2$; see
Soifer and the covering surveys cited below.

The argument is a finite case analysis on ownership of the centre and of the
four outer midpoints, followed by stepwise normalisation / reclassification
and the elimination of joining modes.  Algebraic exclusions are discharged by
exact rational identities, positive Bernstein certificates on compact boxes,
interval certificates, and SMT UNSAT queries; selected geometric interfaces
are sketched in a Lean~4 skeleton.  Machine certificates and the Lean skeleton
(no \texttt{sorry}, no extra axioms \emph{in that skeleton}) support the
narrative; Boolean status flags in Lean are human-maintained metadata, not
geometric theorems.  We do not claim independent peer review, nor a confirmed
resolution of the conjecture, nor a full Lean re-proof of every metric
inequality in the long verification manuscript.
\end{abstract}

\subjclass[2020]{Primary 52C15; Secondary 52C17, 68Q60, 68V15}
\keywords{unit square covering, $S(n)$, $\Pi(n)$, computer-assisted proof,
Bernstein certificates, SMT, Lean~4}

\maketitle

\tableofcontents

\section*{Guide for referees}
This short paper is a \emph{candidate} computer-assisted argument, not a
peer-reviewed resolution.  Read in this order: (i)~Definitions and
Proposition~\ref{prop:M1} (normalisation honesty); (ii)~Proposition~\ref{prop:M3}
and Appendix~\ref{app:M3} (branch tables); (iii)~Propositions in
Sections~\ref{sec:early}--\ref{sec:finale}, each stating
\emph{hypothesis~$\Rightarrow$ conclusion} and citing long-proof modules by
ID; (iv)~Appendix~\ref{app:M1}--\ref{app:M4} (closure patches M1--M4);
(v)~repository \texttt{referee\_pack/REFEREE\_BRIEF.md} for a 30-minute /
2-hour / full-day checklist.  Boolean Lean flags are status metadata.
Do \emph{not} treat dual-solver UNSAT alone as a machine-independent proof.

%% ============================================================
\section{Introduction}\label{sec:intro}
%% ============================================================

Covering a large square by equal smaller squares is a classical theme in
discrete and metric geometry.  Write
\begin{equation}\label{eq:Sn}
S(n)
\;=\;
\sup\bigl\{\,s>0 :
\text{a closed square of side~$s$ admits a covering by~$n$ closed unit squares}
\,\bigr\},
\end{equation}
where each unit square may be translated and rotated and overlaps are allowed.
Soifer's function $\Pi(n)$ is the least number $m$ of unit squares for which
there \emph{exists} some $\varepsilon>0$ such that a square of side
$n+\varepsilon$ admits a covering by $m$ unit squares
(equivalently: there exists $\varepsilon_0>0$ such that every side in
$(n,n+\varepsilon_0]$ admits such a covering);
see~\cite{Soifer2006,SoiferCoverUp}.  With this existential reading, and with
the classical construction that seven unit squares cover some square of side
strictly greater than~$2$ (hence $\Pi(2)\le 7$), one has
\begin{equation}\label{eq:equiv}
S(6)=2
\quad\Longleftrightarrow\quad
\Pi(2)=7.
\end{equation}
The lower bound $S(6)\ge 2$ is elementary: four axis-aligned unit squares tile
$[0,2]^2$.  The candidate upper bound $S(6)\le 2$ is the content of the present
paper; together with $\Pi(2)\le 7$ it would yield $\Pi(2)=7$.

Januszewski proved that covering a square of side $2+\varepsilon$ requires at
least six unit squares, i.e.\ $\Pi(2)\ge 6$ \cite{Januszewski2009}.
Exact small values such as $S(5)=2$ appear in the literature of
Friedman--Paterson and Januszewski; see also Friedman's covering surveys
\cite{FriedmanPaterson2006,FriedmanCovering}.
In January~2026, D\'osa, L\'angi and Tuza revisited congruent square coverings,
determined several small $S(n)$ and boundary-covering variants, and stated as
Conjecture~1.1 that $S(6)=2$ (equivalently $\Pi(2)=7$) \cite{DLT2026}.

\begin{theorem}[Main claim — candidate]\label{thm:main}
Candidate statement: $S(6)=2$.  Equivalently (given $\Pi(2)\le 7$ from known
seven-square constructions), $\Pi(2)=7$: six closed unit squares cannot cover
any square of side strictly greater than~$2$, while a square of side~$2$ admits
such a covering.  This package presents a computer-assisted candidate argument;
it is under revision after a 2026-10-21 public-snapshot review and is
\emph{not} asserted as independently confirmed.
\end{theorem}

The proof occupies Sections~\ref{sec:defs}--\ref{sec:finale}.
Section~\ref{sec:certs} records the logical r\^oles of the computational
certificates.  Related work is collected in Section~\ref{sec:related}.

\begin{remark}[Accuracy tier]
Theorem~\ref{thm:main} is offered as a \emph{computer-assisted candidate}: an
internal audit slice (round~65), certificate inventories, and a Lean~4
obligation skeleton with human-maintained Boolean status metadata for
obligations~1--4.  Lean ``no \texttt{sorry}'' applies to that skeleton only;
Boolean flags are \emph{not} geometric theorems.  This does \emph{not} assert
independent peer review, mathematical confirmation of $S(6)=2$, or a full Lean
re-proof of every metric inequality in the long manuscript.  See the repository
\texttt{DISCLOSURE.md} and \texttt{CHANGELOG\_zh.md}.
\end{remark}

%% ============================================================
\section{Definitions and standing notation}\label{sec:defs}
%% ============================================================

\begin{definition}
For each integer $n\ge 1$, $S(n)$ is defined by~\eqref{eq:Sn}.
Soifer's $\Pi(n)$ is the least integer $m$ for which there \emph{exists} some
$\varepsilon>0$ such that a square of side $n+\varepsilon$ admits a covering by
$m$ unit squares.  Equivalently, there exists $\varepsilon_0>0$ such that every
square of side in $(n,n+\varepsilon_0]$ admits such a covering.
(The universal quantifier ``for every $\varepsilon>0$'' is incorrect: for fixed
$m$, large $\varepsilon$ violate the area bound $(n+\varepsilon)^2\le m$.)
\end{definition}

Throughout the upper-bound argument we work with the centred target
\begin{equation}\label{eq:M1}
T \;=\; [-\mu,\mu]^2,\qquad \mu>1,
\end{equation}
and derive a contradiction from the assumption that six closed unit squares
cover~$T$.  (The translate $[0,2\mu]^2$ is equivalent for all length budgets.)

Write $O$ for the centre of~$T$, label the four corners $1,2,3,4$
counter-clockwise from the lower left, and write $\BOTTOM$, $\RIGHT$, $\TOP$,
$\LEFT$ for the outer edges.  An \emph{outer mid} is the midpoint of an outer
edge.  The central horizontal and vertical target segments are denoted $H$ and
$V$.  A \emph{true gap} on an outer edge is a nonempty open subinterval of that
edge not covered by the two adjacent corner tiles; endpoints of gaps are never
free hypotheses and must be recovered from openness, convexity and closedness.

``Corner tiles abut'' on an outer edge means that the two adjacent corner tiles
together cover the entire edge (overlaps allowed).  Otherwise the uncovered
open interval is a true gap.

%% ============================================================
\section{Roles, normalisation, and parallel budgets}\label{sec:basics}
%% ============================================================

\begin{lemma}[Roles and midpoints]\label{lem:roles}
Four distinct tiles $V_1,V_2,V_3,V_4$ cover the four corners of~$T$; a fifth
tile $C$ covers~$O$; the sixth tile is denoted~$F$.  Each unit tile contains at
most one outer mid.  The role selection does not forbid $F$ from containing a
corner, nor a mid from lying in both a corner tile and a non-corner tile.
\end{lemma}

\begin{proof}
Distinct target corners are at distance at least $2\mu>\sqrt2$, hence cannot
lie in one unit tile.  The distance from a corner to the centre is
$\sqrt2\,\mu>\sqrt2$, so $C$ is not a corner tile.  Distinct outer mids are at
least $\sqrt2\,\mu>\sqrt2$ apart.
\end{proof}

\begin{lemma}[Dominated normalisation]\label{lem:norm}
Fix a target corner as origin with inward coordinates $X,Y\ge 0$.  The unit
tile containing that corner may be enlarged inside~$T$ to the form
\begin{equation}\label{eq:M2}
0\le cX+sY\le 1,\qquad
-\beta\le -sX+cY\le 1-\beta,
\quad c,s\ge 0,\ c^2+s^2=1,
\end{equation}
without moving any other tile and without deleting any previously covered
point of~$T$.
\end{lemma}

\begin{proof}
Among the four oriented unit normals choose one in the \emph{closed} first
quadrant (so $c,s\ge 0$; axis-aligned tiles have a normal on the boundary of
the open first quadrant).  Its original projection interval
$[\alpha,\alpha+1]$ contains~$0$, hence $\alpha\le 0$.  Replacing it by
$[0,1]$ does not remove any point of~$T$ with nonnegative projection.  Keep
the orthogonal unit interval unchanged.
\end{proof}

\begin{proposition}[Dominated corner enlargement; stepwise reclassification]\label{prop:M1}
Starting from any covering of~$T$, apply Lemma~\ref{lem:norm} to a corner tile
$V_i$, leaving the other tiles fixed.  The resulting covering still covers~$T$
and does not delete previously covered points of~$T$.  The same push applied a
second time to that same tile is the identity on the already enlarged tile
(idempotence of that single operator).

This does \emph{not} imply that every later local normalisation or clamping
step (for example, the $\beta\le s/c$ clamp used for anchored corner tiles in
the long verification manuscript) acts as the identity, nor that classification
flags computed before an enlargement remain valid afterwards.  After each
enlargement that may create duplicate hosts or fill gaps, the covering must be
\emph{reclassified} on the new configuration; see the round~65 audit notes on
$\beta$-clamping and stepwise reclassification in the ancillary package
(\texttt{artifacts/still\_S/S6\_round65\_full\_audit.md}, §4).
\end{proposition}

\begin{lemma}[Two nonempty parallel chords]\label{lem:B}
If a unit tile meets two parallel lines at distance $d>1$, then $d\le\sqrt2$
and the sum of the two chord lengths is at most
\begin{equation}\label{eq:M4}
B(d)=\begin{cases}
d-\sqrt{d^2-1},& 1<d\le 3/(2\sqrt2),\\
2(\sqrt2-d),& 3/(2\sqrt2)\le d\le\sqrt2.
\end{cases}
\end{equation}
In particular $B(d)<1$.  Tangency counts as nonempty; an empty chord may not
be substituted into the bound.
\end{lemma}

\begin{proof}[Proof sketch]
With direction parameters $c\ge s>0$ the transverse width is $w=c+s$.  The two
vertical lines lie at $t$ and $t+d$ with $0\le t\le w-d$.  Since $d>1\ge c$,
both chords fall in the triangular end-caps and the total length equals
$(w-d)/(cs)=2(w-d)/(w^2-1)$.  Maximising over $d\le w\le\sqrt2$ yields the
two branches of~\eqref{eq:M4}.  The Euclidean chord-length identities matching
the triangular-cap formulae are recorded in the Lean development of the
parallel-budget geometry (Obligation~1).
\end{proof}

Further elementary budgets (grid contribution at most~$2$ per corner tile;
$L+S+U<2$ for a normalised corner tile; triangular area at most $1/2$) are
collected as Lemmas~R2--R3 in the long verification manuscript and used below
without restating the proofs.

%% ============================================================
\section{Top-level classification}\label{sec:class}
%% ============================================================

If $F$ also contains~$O$, call the configuration class~G5.  Otherwise set
\begin{equation}\label{eq:M5}
k \;=\; \#\{\text{outer mids belonging to }C\cup F\}\in\{0,1,2\}.
\end{equation}
Then $k=0$ is G1, $k=1$ is G2, and $k=2$ with adjacent (resp.\ opposite) mids
is G4 (resp.\ G3).  Each mid not yet supplied by $C\cup F$ chooses an adjacent
corner host injectively.  Reflection reduces host representatives.

\begin{proposition}[Exhaustive branch table]\label{prop:M3}
Every covering of~$T$ falls into exactly one of the classes
G5, G1, centre-owning G2, G4, G3, or residual G2
($k=1$ with the unique mid owned by~$F$ on~$\BOTTOM$).
The residual G2 splits into a near-end host representative
$(\RIGHT,\TOP,\LEFT)=(2,3,4)$ and a far-end representative $(2,4,1)$.
\end{proposition}

The global elimination order is: G5 for all three values of~$k$; then G1,
centre-owning G2, G4 and G3; then near-end residual G2; then reduction of the
far-end representative to the final standard form; then elimination of all
joining modes of that form.

\begin{remark}[Branch tables]
Explicit tables (M3.0--M3.3) and the Boolean near-end checker appear in
Appendix~\ref{app:M3} and \texttt{artifacts/closure/}.  The short-paper
propositions below state hypothesis~$\Rightarrow$ conclusion and cite
long-proof module IDs; they are not self-contained geometric proofs of every
leaf.
\end{remark}

%% ============================================================
\section{Elimination of G5, G1, centre-G2, G3 and G4}\label{sec:early}
%% ============================================================

\begin{proposition}[G5 impossible]\label{prop:M41}
\emph{Hypothesis.} $F$ also contains~$O$ (class G5), with
$k=\#\{\mathrm{mid}\in C\cup F\}\in\{0,1,2\}$.
\emph{Conclusion.} No such covering exists; consequently in any covering of~$T$
exactly one tile contains~$O$, and $O$ is interior to that tile.
\end{proposition}

\begin{proof}[Proof outline]
Long-proof modules: R2 ($k=2$), R3 ($k=0$), R4--R8 ($k=1$).
For G5, modules R2--R3 eliminate $k=2$ and $k=0$.  For $k=1$, modules
R4--R7 reduce to a triple-abutment configuration with an interior second
centre tile; module~R8 then proves the two-corner inequality
\begin{equation}\label{eq:M7}
L_1+L_2+\min(t_1,t_2)+\tfrac45\max(t_1,t_2)\le\tfrac{14}{5}
\end{equation}
by a Bernstein certificate on forty-five rational rectangles after mapping the
full two-dimensional domain into the unit square, and compares it with the
weighted three-ray bound on the interior tile to reach a contradiction in both
orderings of $(t_1,t_2)$.  Thus G5 is empty.  The remaining five tiles stay at
positive distance from~$O$, so a small neighbourhood of~$O$ is covered by~$C$
alone and $O$ is interior to~$C$.
\end{proof}

\begin{proposition}[G1 impossible]\label{prop:M42}
If the centre tile is unique, then $C\cup F$ contains at least one outer mid.
\end{proposition}

\begin{proof}[Proof outline]
Module~R9 eliminates triple abutment by a quarter-boundary budget and reduces
to opposite or adjacent residual gaps.  For opposite gaps, the corrigendum~R10C
first forces the horizontal chord of~$F$ to be strictly positive (via
$g_1+h_C>1$ against $B(\mu)<1$), then combines parallel-chord budgets of $C$
and~$F$ with the corner contribution bound~$<8$.  Adjacent gaps are eliminated
in R10~\S\S3--6 by real endpoints and direction constraints.
\end{proof}

\begin{proposition}[Centre cannot own the unique mid]\label{prop:M43}
If $k=1$, the unique mid belonging to $C\cup F$ cannot belong to~$C$.
\end{proposition}

\begin{proof}[Proof outline]
Axial propagation (R11) yields necessary outer contacts.  Module~R12 treats
the two host representatives when only~$C$ owns the mid, covering all true-gap
and triple-abutment patterns on $\LEFT$, $\TOP$ and~$\RIGHT$.  Shared ownership
$C\cap F$ of the same mid is eliminated separately in R12~\S10 by vertical
redundancy and a true-gap triangle constraint.
\end{proof}

\begin{proposition}[G4 impossible]\label{prop:M51}
No covering has $C$ and~$F$ owning distinct adjacent outer mids.
\end{proposition}

\begin{proof}[Proof outline]
Normalise so that $C$ owns the left mid and $F$ owns the bottom mid.  The three
injective host representatives $(\TOP,\RIGHT)\in\{(4,3),(4,2),(3,2)\}$ are
eliminated in modules R20--R23 by horizontal-sign analysis, forced suspension
segments, and projection budgets on~$C$.  Shared left-mid ownership by the free
corner tile is ruled out before the last representative.
\end{proof}

\begin{proposition}[G3 impossible]\label{prop:M52}
No covering has $C$ and~$F$ owning distinct opposite outer mids.
\end{proposition}

\begin{proof}[Proof outline]
Normalise so that $C$ owns the left mid and $F$ owns the right mid.  Three host
representatives (after reflection) are eliminated in R24--R25 by abutment/gap
analysis on $\TOP$ and~$\BOTTOM$, empty/nonempty vertical chords of~$F$, and a
local exit lemma for nonempty cases.
\end{proof}

\begin{corollary}[Universal centre forbids outer mids]\label{cor:M53}
In any covering of~$T$ the unique centre tile contains no outer mid, and the
remaining non-corner tile contains some outer mid.
\end{corollary}

\begin{proof}
Combine Propositions~\ref{prop:M41}--\ref{prop:M52}.
\end{proof}

%% ============================================================
\section{Near-end residual G2}\label{sec:near}
%% ============================================================

After Corollary~\ref{cor:M53} the only surviving class is residual G2 with the
unique mid owned by~$F$ on~$\BOTTOM$.

\begin{proposition}[Near-end impossible]\label{prop:M61}
\emph{Hypothesis.} Residual G2 with $F$ owning the bottom mid, $C$ owning no
outer mid, and host injection $(\RIGHT,\TOP,\LEFT)=(2,3,4)$ (free corner
$W=V_1$).
\emph{Conclusion.} No such covering of~$T$ exists.
\end{proposition}

\begin{proof}[Proof outline]
See long-proof modules R13--R19 and R26--R28, summarised by the Boolean
branch table in Appendix~\ref{app:M3} (Table~M3.1).  Module~R18 shows that
either both $\RIGHT$ and~$\TOP$ abut, or exactly one of them carries a true
gap owned by~$C$.  In the double-abutment case: R13 eliminates nonempty
horizontal chords of~$F$; R14--R15 force left abutment when that chord is
empty and produce an interior point $Q_0$ of~$C$; R15S/R16P supply the
strict angular separation and height hypotheses; R16 then covers the
necessary continuous domain by $1173$ rational boxes with interval
certificates excluding every leaf.  Remaining left-gap / nonempty-$F$
patterns are eliminated in R17--R18.  Top-gap patterns fall to R19 and~R26;
right-gap patterns fall to R27--R28 (both rotation signs).  A normalisation
that fills a gap returns to an already eliminated abutment leaf
(Appendix~\ref{app:M1}).
\end{proof}

%% ============================================================
\section{Far-end reduction to standard form}\label{sec:far}
%% ============================================================

Henceforth fix the far-end representative: $V_1$ owns the left mid, $V_2$ the
right mid, $V_4$ the top mid, and the free corner tile is $W=V_3$.  The centre
tile $C$ owns no outer mid; $F$ owns the bottom mid and not the centre.

\begin{proposition}[Far-end standard form]\label{prop:M71}
\emph{Hypothesis.} Residual G2 with $F$ owning the bottom mid, $C$ owning no
outer mid, and far-end hosts $V_1\supset\mathrm{LEFT\,mid}$,
$V_2\supset\mathrm{RIGHT\,mid}$, $V_4\supset\mathrm{TOP\,mid}$, free corner
$W=V_3$.
\emph{Conclusion.} The covering satisfies: $\LEFT$, $\RIGHT$ and $\TOP$ abut
via corner tiles; $F\subset\{y<0\}$; $C$ avoids~$\BOTTOM$; and one copy of~$F$
contains both closed endpoints of the bottom gap together with the lower axial
tip of~$C$ (the \emph{final standard form}).
\end{proposition}

\begin{proof}[Proof outline]
See Table~M3.2 in Appendix~\ref{app:M3} and frozen module R29N in
Appendix~\ref{app:M2}.

\emph{Step~1.} Modules R27 and~R29 eliminate standardised top true gaps for
both rotation signs of~$C$.  Thus $\TOP$ abuts.  Any lateral true gap lies too
far from the bottom mid for~$F$ to reach, so lateral gaps are owned entirely
by~$C$.  Since $C$ cannot meet both $\LEFT$ and~$\RIGHT$ at distance
$2\mu>\sqrt2$, the only remaining cases are left gap, right gap, or triple
abutment (Proposition~R29N.1).

\emph{Step~2.} If $F\cap H\neq\varnothing$, Proposition~R29N.2 forces $C$ to
avoid~$\BOTTOM$; R31 eliminates triple abutment and R32--R36 eliminate the
lateral-gap cases.  Hence $F\cap H=\varnothing$.  If some point of~$F$ had
$y\ge 0$, convexity with $(0,-\mu)\in F$ would produce a point of $F\cap H$
inside~$T$, a contradiction.  Thus $F\subset\{y<0\}$.

\emph{Step~3.} Corollary~\ref{cor:M53} applied to an inscribed target yields
axial-root caps $\le 2-\mu$ (module~R37) without circularity.  Modules
R37--R39 eliminate all right and left true-gap patterns (both rotation signs,
direct and bridged contacts).

\emph{Step~4.} In the residual triple-abutment case, R39 gives the upper axial
root $u>38/79$, whence $C$ avoids~$\BOTTOM$.  The three anchored corner tiles
can contribute only their matching central half-axes, so the entire bottom gap
and the lower tip of~$C$ are forced into the same tile~$F$.
\end{proof}

%% ============================================================
\section{Standard-form notation and four joining modes}\label{sec:std}
%% ============================================================

Let $a,x,z\in(0,1)$ be the half-angle parameters of $V_1,V_2,V_4$, and let
$h\in[0,\sqrt2-1]$ be the half-angle of the smaller axial angle of~$C$.  Set
\begin{equation}\label{eq:M8}
b(t)=\frac{1+t^2}{1-t^2},\quad
k(t)=\frac{2t}{1-t^2},\quad
g(t)=\frac{1-t}{1+t},\quad
q(t)=\frac{t}{2-t},
\end{equation}
and write $\delta(t)=b(t)-1$, $\Delta_z=\delta(z)$.  The four actual axial tips
of~$C$ are $L=(-l,0)$, $R=(r,0)$, $T_C=(0,u)$ and $B_C=(0,-d)$.  Writing
$S_1,S_2$ for the actual short legs of the bottom corner tiles and
$\Sigma=S_1+S_2$, Proposition~\ref{prop:M71} yields
\begin{equation}\label{eq:M9}
\Sigma>2-\sqrt2>\tfrac7{12},\qquad
\Sigma>(1-d)^2,\qquad
\Sigma>2-\frac1{1-d}.
\end{equation}
Left abutment supplies the constraints
\begin{equation}\label{eq:M10}
0<z<a<1,\quad a+2z<1,\quad z<\tfrac4{15},\quad
0\le S_1<\mathcal F(a,z):=b(a)-k(a)\bigl(2-g(z)\bigr).
\end{equation}

\begin{lemma}[Normalised seam point in~$C$]\label{lem:M81}
There is a point
\begin{equation}\label{eq:M11}
P_1=(-1+X,e)\in C,\qquad
X=\frac{b(a)+S_4-2}{k(a)-k(z)},\quad
e=1-S_4-k(z)X,
\end{equation}
obtained from the real-scale seam $P_\mu\in C$ by a convexity argument inside
the same tile~$C$ (modules R40 and~R42).  The construction does not rescale the
covering to $\mu=1$ and does not move~$C$.
\end{lemma}

Let $Q_2$ and $Q_U$ be the actual inward tips of $V_2$ and $V_4$ on the positive
axes.  Whether $C$ contains these tips yields four exhaustive joining modes:
double-direct, upper bridge by~$W$, right bridge by~$W$, and double bridge.
If $C$ misses a tip, the corresponding open axial gap can be filled only by the
same free corner tile~$W$, which is then forced by closedness to contain both
the tip of~$C$ and the corner tip.

%% ============================================================
\section{Elimination of upper bridges and negative rotations}\label{sec:M9}
%% ============================================================

\begin{proposition}\label{prop:M91}
In the standard form of Proposition~\ref{prop:M71}, the free corner tile~$W$
cannot contain~$Q_U$.
\end{proposition}

\begin{proof}[Proof outline]
Modules R44--R45 show that an upper bridge forces $z\le\tau<2/5$ for the
parameter $\tau$ of~$W$ and empties $W\cap H$, whence $C$ contains~$Q_2$
directly and double bridging is impossible.  Modules R46 and~R48 produce a
real exit point $T_R\in C\cap W\cap V_2$.  Small angles $h\le 1/24$ fall to~R52.
Positive rotations are eliminated in R53--R55 via the seam-derived bound
$S_1<q(Q)$ and the two normals of~$C$ at~$T_R$.  Negative rotations are reduced
in R56--R58 to $0<\tau<1/6$ and then excluded in~R59 by the area-forced
negativity
\begin{equation}\label{eq:M12}
G=H_F-q(x+\delta)-g(x)+k(x)(B_t-1)<0,
\quad H_F=2-\frac1{1-\delta-\tau},
\end{equation}
using endpoint polynomial certificates on small/medium~$x$ and a left-root
projection contradiction on large~$x$.
\end{proof}

Consequently only double-direct and right-bridge modes remain, both with
$C\ni Q_U$ and hence $u>1-z$.  In the right-bridge case the horizontal chord
length of~$W$ is strictly less than~$\Delta_z$, which after establishing that
the relevant parameters are~$<1$ yields
\begin{equation}\label{eq:M13}
S_2<q(r+\Delta_z);
\end{equation}
double-direct admits the stronger $S_2<q(r)$.

\begin{proposition}\label{prop:M92}
All negative rotations in the two residual joining modes are impossible.
Every surviving positive rotation satisfies $1/10<z<1/5$.
\end{proposition}

\begin{proof}[Proof outline]
Small-angle double-direct and right-bridge cases are eliminated in R43 and
R49--R51 respectively, both forcing $h>1/8$.  Module~R60 partitions the
negative-rotation domain into four regions and contradicts~\eqref{eq:M9} in
each, paying the loss~\eqref{eq:M13}.  Modules R61--R62 eliminate
$z\ge 1/5$ and $z\le 1/10$ for positive rotations by asymmetric triangles on
the same~$F$ and concavity of the budget form, leaving the open strip
$1/10<z<1/5$.
\end{proof}

%% ============================================================
\section{Right bridges are impossible}\label{sec:M10}
%% ============================================================

\begin{proposition}\label{prop:M101}
In the residual domain of Proposition~\ref{prop:M92}, $W$ does not contain the
right axial tip~$R$ of~$C$.  Consequently the entire right-bridge class is
empty.
\end{proposition}

\begin{proof}[Proof outline]
Module~R63 first obtains $r<17/20$ from axial chords and the $R/T_C$
projection.  Assuming $W\ni R$, the tile meets both $\TOP$ and~$H$, forcing a
top leg $D>1-\Delta_z>11/12$ and a half-angle
\begin{equation}\label{eq:M14}
t<\tfrac18,\qquad t<\tfrac97\Delta_z.
\end{equation}
Comparing real projections in both orientations of~$W$ yields an abscissa
$\xi$ of the contact on~$H$ satisfying
\begin{equation}\label{eq:M15}
\xi<\Delta_z+t
\quad\text{or}\quad
\xi>1-t>\tfrac78.
\end{equation}
The second alternative contradicts $r<17/20$ if the contact is~$R$.  The first,
combined with~\eqref{eq:M13}--\eqref{eq:M14} and left long-leg loss, produces
\begin{equation}
\Sigma<\Phi\!\Bigl(\frac{2z}{1+z}\Bigr)+q\!\Bigl(\frac{23}{7}\Delta_z\Bigr),
\quad
\Phi(w)=(1+w)\bigl(1-\sqrt{w(2+w)}\bigr),
\end{equation}
which is decreasing on $1/10\le z\le 1/5$ and strictly less than $7/12$ at the
left endpoint, contradicting~\eqref{eq:M9}.
\end{proof}

%% ============================================================
\section{Double-direct positive rotations}\label{sec:finale}
%% ============================================================

\begin{proposition}\label{prop:M111}
No positive-rotation double-direct covering survives in the residual strip of
Proposition~\ref{prop:M92}.
\end{proposition}

\begin{proof}[Proof outline]
From~\eqref{eq:M10} and $1/10<z<1/5$ one obtains $2/7<a<3/5$.  Double-direct
gives $S_2<q(r)$.  At the formal endpoint $S_4=g(z)$ define
\begin{equation}\label{eq:M16}
X_G=\frac{(1-z)\bigl[a^2(1+2z)-z\bigr]}{(a-z)(1+az)},\qquad
e_G=\frac{2az(1-a-z)}{(a-z)(1+az)}.
\end{equation}
After verifying $a<g(h)$ one has the projection bound
\begin{equation}\label{eq:M17}
r<R_P(h):=\delta(h)+X_G+k(h)e_G,
\end{equation}
while the pair $(T_C,R)$ yields $r<R_T(h):=h+z/k(h)$ for $h>0$.

On $0\le h\le 1/4$ one has $a<3/5\le g(h)$; monotonicity of~\eqref{eq:M17}
reduces the right-root upper bound to $h=1/4$, after which
$\mathcal F+q\bigl(R_P(1/4)\bigr)<7/12$ on the full $(a,z)$-rectangle.
On $1/4\le h\le 1/3$, concavity of the real $P_1/T_C$ relation forces $a>2/5$,
and four endpoint values control $\Sigma<7/12$.  On
$1/3\le h\le\sqrt2-1$, a rational outer cap $29/70$ splits according as
$a\ge g(h)$ or not; both subcases again force $\Sigma<7/12$.

The fourteen deduplicated rational targets are certified by twenty-eight
numerator/denominator polynomials whose one hundred fourteen strictly positive
rational Bernstein coefficients, together with reconstruction identities, cover
the whole continuum $[0,\sqrt2-1]$ (module~R64).  Every angular segment
contradicts~\eqref{eq:M9}.
\end{proof}

\begin{proof}[Proof of Theorem~\ref{thm:main}]
Assume six unit squares cover~$T$ with $\mu>1$.  Proposition~\ref{prop:M3}
places the covering in one of the top-level classes.  Propositions
\ref{prop:M41}--\ref{prop:M52} eliminate G5, G1, centre-owning G2, G3 and~G4.
Proposition~\ref{prop:M61} eliminates the near-end residual G2.
Proposition~\ref{prop:M71} reduces the far-end residual G2 to the final
standard form.  Propositions~\ref{prop:M91}--\ref{prop:M111} eliminate every
joining mode, rotation sign and residual parameter domain of that form.
Hence (subject to the completeness of the cited modules and certificates)
$\mu>1$ is impossible under the candidate narrative.  Combined with the
elementary covering of side~$2$ and the classical $\Pi(2)\le 7$, this would
yield $S(6)=2$ and therefore $\Pi(2)=7$ by~\eqref{eq:equiv}.  The public
package remains a candidate under revision.
\end{proof}

%% ============================================================
\section{Computational certificates}\label{sec:certs}

The geometric narrative of Sections~\ref{sec:basics}--\ref{sec:finale} is
supported by three layers of evidence.

\paragraph{Exact rational certificates.}
On a polynomial over $[0,1]^m$, the tensor Bernstein basis is nonnegative and
sums to one.  Nonnegative (resp.\ strictly positive) coefficients prove
nonnegativity (resp.\ strict positivity).  Parameter domains are mapped into
rectangles by explicit rational substitutions after denominators are shown
positive.  Verifiers regenerate targets, rebuild basis-change identities and
rational coefficients, rather than trusting a generator's claimed minimal
coefficient.  Module~R8 additionally checks a complete prefix-tree cover;
module~R64 checks every polynomial on its fixed rectangles.

\paragraph{Continuous interval certificates.}
The R16 root box outwardly contains the necessary domain.  Each internal node
splits into children whose union recovers the parent; each leaf is excluded by
exact fractional envelopes.  Denominator positivity, monotonic substitution
direction, and all leaves together with their boundaries are checked.

\paragraph{Real-algebraic UNSAT queries.}
Selected propositions are accompanied by QF\_LRA or QF\_NRA queries.  Only
UNSAT is accepted.  Queries are classified as (E)~exact certificates,
(H)~hand proofs with solver cross-check, or (S)~dual-solver UNSAT (z3 and
cvc5) not yet reduced to~(E).  Dual-wrapper spot checks on the M8--M11
inventory report pass on the expanded faithfulness sample shipped with the
ancillary package.  Finite floating-point sampling is never used as an
exclusion certificate.

\paragraph{Lean~4 fragments.}
Selected interfaces---parallel-budget geometry and Euclidean chord lengths;
partial ownership / role bridges; M7.1 branch-logic scaffolding; far-end PDF
faithfulness sketches; and algebraic pieces used in M8--M11 together with
near-end leaf wirings---are present in the Lean~4 package \texttt{S6Core}
(mathlib).  The shipped skeleton builds with no \texttt{sorry} and no extra
axioms \emph{as a skeleton}.  Boolean ``closed'' flags
(e.g.\ \texttt{wholeTheorem\_S6\_eq\_2\_closed}) are human-maintained
\emph{status metadata}, not theorems about six unit squares covering a target.
Several leaf conclusions are still propositional placeholders of the form
\texttt{living $\to$ True}; ownership-bridge classifiers are not yet proved
faithful to the manuscript's $k$ / $F\ni O$ / sixth-tile role selection.
The long verification manuscript remains the geometric reference for metric
inequalities discharged primarily by certificates.

\paragraph{Reproduction.}
The public repository ships Lean sources under \texttt{lean/}, certificate
registries and spot-check checkers under \texttt{artifacts/}, and the long
verification manuscript under \texttt{artifacts/long\_proof/}
(\texttt{S6\_full\_proof.md}/\texttt{.pdf}).  Companion certificate JSON files
needed by the shipped R42/R54--55/R64 checkers sit next to those scripts.
A compressed copy of the historical \texttt{ex/all/} verification tree is
shipped as \texttt{artifacts/ex\_all.tar.gz} with SHA-256 sums in
\texttt{artifacts/SHA256SUMS.txt}; an unpacked copy sits at
\texttt{artifacts/ex/all/}.  Closure patches M1--M4 live in
\texttt{artifacts/closure/} and in Appendices~\ref{app:M1}--\ref{app:M4}.
See \texttt{referee\_pack/HOW\_TO\_REPRODUCE.md}.  A minimal path from a clean
clone is
\begin{verbatim}
cd artifacts/checkers && python3 check_round64.py && python3 check_round63.py
cd ../closure && python3 m3_branch_checker.py
# optional: cd ../../lean/S6Core && lake build
\end{verbatim}

%% ============================================================
\section{Related work}\label{sec:related}
%% ============================================================

Soifer's JCTA note and Cover-Up Squared writings introduced $\Pi(n)$
\cite{Soifer2006,SoiferCoverUp}.
Januszewski's Monthly note established $\Pi(2)\ge 6$ \cite{Januszewski2009}.
Friedman--Paterson and Friedman's surveys record constructions and small-case
discussions \cite{FriedmanPaterson2006,FriedmanCovering}.
D\'osa--L\'angi--Tuza proved several exact small values and boundary analogues
and isolated the conjecture $S(6)=2$ \cite{DLT2026}; Theorem~\ref{thm:main}
is a candidate argument addressing that conjecture.
Asymptotic packing-waste problems of Chung--Graham
\cite{ChungGraham2009,ChungGraham2020} concern a different regime.

On computer-assisted proofs in geometry, we follow the AMS guidelines on
independent checkability and durable availability of code and data
\cite{AMSComputerAssisted}.

%% ============================================================
\section{Concluding remarks}\label{sec:conclusion}
%% ============================================================

If confirmed, Theorem~\ref{thm:main} would determine the first previously open
value of $S(n)$ beyond the classical small cases.  The same case-analysis and certificate
pipeline suggests approaches to neighbouring values (e.g.\ refinements of
$\Pi(3)$ and boundary-covering variants $S_{\mathrm{bd}}(n)$ from~\cite{DLT2026}),
though those questions lie outside the present paper.

\section*{Acknowledgements}
This manuscript and the accompanying certificate / Lean scaffolding were
produced by a generative AI system and released publicly by Ningyu Zhou
(School of Physics, Peking University).  The releaser hosts and maintains the
package and does not claim independent mastery of every step.  Independent
external peer review has not occurred.  See the repository
\texttt{DISCLOSURE.md} for the full release ethics statement.


%% ============================================================
\appendix
\section{Closure patch M1: global normalisation honesty}\label{app:M1}
%% ============================================================

The maturity revision weakened Proposition~\ref{prop:M1}: a single dominated
enlargement is idempotent, but later $\beta$-clamps need not be the identity,
and classification must be recomputed after enlargements that fill gaps.
The following propositions, adapted from the Chinese closure note
\texttt{artifacts/closure/M1\_proposition\_zh.md}, record the stronger
\emph{intended} global protocol used by the long proof.  They are packaging
and logical-cycle closures, not new geometric exclusions.

\begin{proposition}[One-shot global corner normalisation]\label{prop:M11}
Apply Lemma~\ref{lem:norm} once to each corner tile $V_1,\dots,V_4$, leaving
$C$ and $F$ fixed, and (when $c,s>0$) clamp the orthogonal parameter $\beta$
into the interval determined by $s/c$ and $1-c/s$ as in the long manuscript.
Then: coverage of~$T$ is nondecreasing; roles of $C$ and $F$ are unchanged
up to non-increasing $k$; open-first-quadrant normals are unique so the
procedure is canonical; and no corner tile acquires $O$ or two distinct outer
mids.
\end{proposition}

\begin{proposition}[Classify once after global normalisation]\label{prop:M12}
On the output of Proposition~\ref{prop:M11}, perform the G1--G5 / host /
near--far / abutment--gap / $F\cap H$ classification of
Section~\ref{sec:class} exactly once.  Subsequent local ``normalisation''
steps of types (L$\alpha$) and (L$\beta$) on already-normalised corner tiles
are the identity, so there is no ``normalise $\to$ reclassify $\to$
normalise'' loop inside elimination modules.
\end{proposition}

\begin{remark}
Proposition~\ref{prop:M12} assumes no module requires a strictly
\emph{non}-normalised corner tile.  The module-by-module check is recorded in
\texttt{artifacts/closure/REVIEW\_addendum\_item2\_zh.md}.  If a future module
needs a non-normalised intersection, that module must be rewritten or the
global protocol delayed with a well-founded order.  Until independently
verified, treat Propositions~\ref{prop:M11}--\ref{prop:M12} as candidate
closure text accompanying the long proof.
\end{remark}

%% ============================================================
\section{Closure patch M2: frozen module R29N}\label{app:M2}
%% ============================================================

References in the short paper and long proof to ``round-29 continuation notes
\S4'' are hereby frozen as module~\textbf{R29N}.  Withdrawn note~\S5 (swapped
parallel directions) is \emph{not} part of the candidate argument.

\begin{proposition}[R29N.1: three outer-edge states]\label{prop:R29N1}
Under the far-end hypothesis of Proposition~\ref{prop:M71} with $\TOP$
abutting, the remaining outer-edge states are exactly: (i)~LEFT gap owned by
$C$, RIGHT and TOP abutting; (ii)~RIGHT gap owned by $C$, LEFT and TOP
abutting; (iii)~LEFT, RIGHT and TOP all abutting.
\end{proposition}

\begin{proposition}[R29N.2: nonempty $F\cap H$ forces $C$ off BOTTOM]\label{prop:R29N2}
If $F\cap H\neq\varnothing$, then $C$ avoids $\BOTTOM$ entirely.
\end{proposition}

\begin{proposition}[R29N.4: triple abutment + nonempty $F\cap H$]\label{prop:R29N4}
Under triple abutment and $F\cap H\neq\varnothing$, the free corner $W$ cannot
contain the top mid.  (Does not migrate to lateral-gap cases.)
\end{proposition}

Full inequality chains, including the 266-coefficient adjacent-corner
certificates retained from R13, are in
\texttt{artifacts/closure/M2\_R29\_module\_zh.md} and the long proof.

%% ============================================================
\section{Closure patch M3: explicit branch tables}\label{app:M3}
%% ============================================================

\begin{center}
\small
\begin{tabular}{@{}llll@{}}
\toprule
Class & Definition & Closing modules & Exhaustion note \\
\midrule
G5 & $F\ni O$; $k\in\{0,1,2\}$ & R2--R8 & Prop.~\ref{prop:M41} \\
G1 & unique centre, $k=0$ & R9--R10C & Prop.~\ref{prop:M42} \\
centre-G2 & $k=1$, mid in $C$ & R11--R12 & Prop.~\ref{prop:M43} \\
G4 & $k=2$, adjacent mids & R20--R23 & Prop.~\ref{prop:M51} \\
G3 & $k=2$, opposite mids & R24--R25 & Prop.~\ref{prop:M52} \\
residual G2 & $k=1$, mid only in $F$ & near Prop.~\ref{prop:M61}; far Prop.~\ref{prop:M71}+M8--M11 & Cor.~\ref{cor:M53} \\
\bottomrule
\end{tabular}
\end{center}
\captionof{table}{Table M3.0 (top-level). Full near-end / far-end / joining
tables: \texttt{artifacts/closure/M3\_branch\_table\_zh.md}.}

\noindent
The Python checker \texttt{artifacts/closure/m3\_branch\_checker.py} verifies
that the near-end Boolean grid (24 cells) is covered by live closing rows
plus explicitly impossible cells.  It is \emph{not} a geometric proof.
Running it on the shipped tree returns \texttt{PASS}.

%% ============================================================
\section{Closure patch M4: E/H/S certificate tiers}\label{app:M4}
%% ============================================================

Computational exclusions in the long proof are tagged:
\begin{itemize}[leftmargin=*]
\item[(E)] exact certificates (Bernstein / Farkas / rational IBB / CAD);
\item[(H)] hand proofs with solver cross-check;
\item[(S)] dual-solver UNSAT (z3 and cvc5) without an independent exact
certificate.
\end{itemize}
On the tracked unique-query inventory accompanying the peer-ready dump,
roughly $60\%$ are (E) or (H) and roughly $40\%$ remain (S); exact counts
drift as Item-3 upgrades land and are recorded in
\texttt{artifacts/closure/M4\_unsat\_inventory.md}.
\textbf{Class (S) must not be rewritten as machine-independent proofs.}
External referees should sample (S) queries for modelling faithfulness even
when both solvers report UNSAT.

\bibliographystyle{amsplain}
\bibliography{refs}

\end{document}
